\section{Graph RAG}
\label{sec:graphrag}

\subsection{Why conventional RAG is insufficient}

Standard retrieval-augmented generation embeds text chunks and retrieves those
nearest a query. This works for lookup questions and fails for three classes of
question that dominate industrial use.

\begin{center}
\renewcommand{\arraystretch}{1.3}
\begin{tabularx}{\linewidth}{@{}X l l@{}}
\toprule
\textbf{Question} & \textbf{Requires} & \textbf{Vector search} \\
\midrule
What does the SOP say about pump lubrication? & lookup & \tick\ works \\
If we isolate P-101A, what is affected downstream? & traversal & \cross\ fails \\
Which vendors supplied valves in corrosion incidents since 2021? & join & \cross\ fails \\
What are the recurring failure themes across five years? & aggregation & \cross\ fails \\
\bottomrule
\end{tabularx}
\end{center}

The failure is structural, not a matter of tuning. Consider the corpus fragment
\code{P-101A --FEEDS--> E-204} and \code{E-204 --FEEDS--> V-1201}. The query
``what is affected if P-101A is isolated'' must return V-1201. But V-1201 never
appears in any sentence containing P-101A, so no embedding of the query is near
any chunk mentioning V-1201. Similarity is not connectivity. Increasing $k$,
changing the embedding model, or reranking cannot fix this; the information
required is in the \emph{topology}, which embeddings discard.

\begin{keybox}{A P\&ID is already a graph}
The strongest argument for this approach in an industrial setting is that
piping and instrumentation diagrams are graphs by construction: equipment are
nodes, process lines are edges. Flattening them into embeddings destroys the
structure the document was drawn to express. Extracting a graph does not
impose a foreign abstraction --- it recovers the native one.
\end{keybox}

\subsection{The ontology}

Constraining extraction to a closed schema is the single largest quality lever.
Free-form extraction produces a graph in which ``Pump 101A'', ``P-101A'' and
``the main feed pump'' are three unrelated nodes, and every traversal returns a
fragment of the truth.

The schema is a configuration artefact, selected by an environment variable:

\begin{lstlisting}[language=yamlx]
packs:
  general:                       # domain-neutral: any corpus
    nodes:
      Person:       [role, affiliation, email]
      Concept:      [definition, domain]
      Technology:   [category, vendor]
      Publication:  [venue, year, doi]
      ...
    relations:
      AUTHORED_BY:  [[Document, Person], [Publication, Person]]
      PART_OF:      [[Section, Document], [Project, Organization]]

  industrial:                    # refinery / plant extension
    nodes:
      Equipment:    [tag, equip_type, unit, description]
      Instrument:   [tag, instr_type, measures, unit]
    tagged: [Equipment, Instrument, Line]
    relations:
      FEEDS:        [[Equipment, Equipment]]
      PART_OF:      [[Equipment, Equipment]]
\end{lstlisting}

Three properties make this robust:

\begin{description}[leftmargin=0pt,style=nextline]
\item[Packs merge rather than overwrite.] \code{PART\_OF} is declared in both
packs. The merged ontology retains every endpoint signature
--- \code{Section$\rightarrow$Document}, \code{Project$\rightarrow$Organization}
\emph{and} \code{Equipment$\rightarrow$Equipment} --- implemented as one
relation table carrying multiple \code{FROM \ldots TO} pairs.
\item[Cross-pack references degrade gracefully.] The industrial pack references
\code{Person}, which lives in the general pack. Enabling industrial alone drops
those signatures instead of failing to start.
\item[Tagging is opt-in.] Only types listed under \code{tagged:} pass through
identifier canonicalisation. A corpus of research papers is never forced
through asset-tag parsing.
\end{description}

\subsection{Entity resolution}

Identifier canonicalisation is where naive implementations quietly fail. The
resolver normalises Unicode, removes stop words, maps spelled-out equipment
nouns to their conventional function codes, and matches an ISA-5.1 shaped
pattern anywhere in the string.

\begin{center}
\renewcommand{\arraystretch}{1.15}
\begin{tabular}{@{}l l@{}}
\toprule
\textbf{Surface form} & \textbf{Canonical key} \\
\midrule
\code{P-101A}, \code{P101A}, \code{PUMP-101A} & \code{Equipment:P-101A} \\
\code{Pump 101-A}, \code{pump P101 A} & \code{Equipment:P-101A} \\
\code{the main pump P-101 A} & \code{Equipment:P-101A} \\
\code{centrifugal pump 101 A}, \code{Pump no. 101 A} & \code{Equipment:P-101A} \\
\midrule
\code{Exchanger 204}, \code{E-204}, \code{heat exchanger E204} & \code{Equipment:E-204} \\
\code{Fire Water} & (not a tag --- name slug) \\
\bottomrule
\end{tabular}
\end{center}

All seven spellings of the first asset collapse to one node. During development
a defect in this component --- a stop-word list that stripped a bare
\code{A} --- silently split \code{P-101A} from \code{P-101}, fragmenting the
graph without any error. It is now covered by regression tests.

\subsection{Retrieval modes}
\label{sec:modes}

Query shape selects the strategy.

\begin{center}
\renewcommand{\arraystretch}{1.3}
\begin{tabularx}{\linewidth}{@{}l X l@{}}
\toprule
\textbf{Mode} & \textbf{Trigger and mechanism} & \textbf{Answers} \\
\midrule
\code{vector} & no structural signal; BM25 + dense fused by RRF & lookup \\
\code{local} & an identifier is present; one-hop neighbourhood + linked chunks & context \\
\code{traverse} & impact vocabulary \emph{and} an identifier; three-hop walk & dependency \\
\code{global} & corpus-wide vocabulary; map-reduce over community summaries & themes \\
\bottomrule
\end{tabularx}
\end{center}

Measured behaviour on the indexed corpus:

\begin{lstlisting}
[traverse] If we isolate P-101A which downstream units are affected?
           -> nodes=7  chunks=3   (V-1201 reached at two hops)
[global  ] what are the recurring failure themes?
           -> communities=2 consulted
[local   ] what does the SOP say about P-101A
           -> nodes=5  chunks=3   (V-1201 NOT present: one hop only)
[vector  ] what is the bearing vibration alarm limit
           -> chunks=3
\end{lstlisting}

The contrast between \code{traverse} and \code{local} is the demonstration: the
same anchor entity, differing only in hop depth, and the indirect dependency
appears only in the deeper walk.

\subsection{Result on a general corpus}

To validate that the ontology is genuinely domain-portable, a corpus containing
three curricula vitae and a fifteen-page electronics journal article was
indexed with \code{ONTOLOGY=general,industrial}.

\begin{center}
\renewcommand{\arraystretch}{1.15}
\begin{tabular}{@{}l r l r@{}}
\toprule
\textbf{Node type} & \textbf{Count} & \textbf{Relation} & \textbf{Count} \\
\midrule
Person & 152 & AUTHORED\_BY & 158 \\
Concept & 94 & MENTIONS & 83 \\
Technology & 68 & HAS\_SKILL & 37 \\
Skill & 35 & USES & 30 \\
Publication & 28 & APPLIED\_TO & 13 \\
Document & 26 & REPORTS\_METRIC & 12 \\
Section & 17 & COMPARED\_WITH & 10 \\
Organization & 14 & AFFILIATED\_WITH & 10 \\
others & 65 & others & 38 \\
\midrule
\textbf{Total} & \textbf{499} & \textbf{Total} & \textbf{391} \\
\bottomrule
\end{tabular}
\end{center}

\noindent Thirty-seven communities were detected and summarised. The industrial
types remained near-empty, as expected for this corpus --- confirming that the
packs neither interfere with one another nor force inappropriate structure onto
the data.

\subsection{A defect worth recording}

The first generation of community summaries stated relationships backwards:
the graph asserted \code{V-4401 PART\_OF P-101A} while the summary read
``P-101A is a component of V-4401''. The cause was that the clustering graph
was built with an undirected \code{networkx.Graph}, so edge orientation was
lost before the summarisation prompt was rendered. Clustering legitimately
requires an undirected view, but rendering does not. The graph is now
constructed as a \code{DiGraph} and converted to undirected only for the Leiden
step. In a system whose central claim is grounded accuracy, confidently
reversed relationships are among the most damaging failures possible, and they
produced no error of any kind.
